\section{Introduction and Clinical Motivation}

Redundant wait times constrain patient throughput in the pediatric Emergency Department (ED) at The Hospital for Sick Children (SickKids). Under standard care, patients queue twice: first for an initial physician assessment to order diagnostics, then again for results. Faster diagnostic turnaround enables earlier treatment decisions, shortens time to disposition, and reduces ED crowding---factors that directly affect patient outcomes and experience.

The hospital has implemented a Machine Learning Medical Directive (MLMD) system that uses triage data to predict and authorize relevant tests proactively, parallelizing diagnostic processing with the initial wait~\cite{singh_assessment_2022}. The MLMD targets four high-volume, time-sensitive diagnostic pathways: arm X-rays, electrocardiograms (ECGs), appendix ultrasounds, and testicular ultrasounds. However, the model's performance across these pathways is limited by significant label noise. Positive labels in the SickKids Electronic Health Record (EHR) are reliable, but negative labels are ambiguous because the EHR is not integrated with external providers. Because records of externally performed tests are unavailable, their occurrence must be inferred retrospectively from the patient's final documented diagnosis. The current baseline system utilizes a regex-driven approach to identify variations of diagnoses in clinical text and map them to their corresponding required tests. While this rules-based mapping is clinician-validated, it is inherently brittle; it struggles to capture the vast semantic nuance, shorthand, and variability inherent to clinical documentation. As a result, many valid diagnoses---and the external tests they imply---are missed by the regex, causing the system to misclassify these instances as true negatives. Across 483,705 triage encounters spanning June 2018 to January 2026, this label ambiguity pushes the MLMD toward a conservative decision boundary that artificially depresses recall. Each relevant test the model fails to authorize returns a patient to the slower, serialized workflow, forfeiting substantial efficiency gains.

To overcome this limitation, we propose a two-pronged methodology to handle asymmetric label noise at both the data and algorithmic levels. First, we replace the brittle regex baseline with a dynamic, evidence-grounded Knowledge Graph and Large Language Model (LLM) pipeline, capable of semantically parsing clinical triage notes to accurately impute missing test labels---recovering 11,110 test-ordering events the regex baseline missed. Second, to manage the residual uncertainty inherent to imputed labels, we apply LiLAW (Lightweight Learnable Adaptive Weighting)~\cite{moturu_lilaw_2025}, a meta-learning framework that dynamically down-weights suspected false negatives during training. By combining semantic knowledge extraction with robust meta-learning, we demonstrate that the MLMD can safely expand its decision boundary, substantially improving recall and unlocking further efficiency gains for the pediatric ED.
